\documentclass[10pt,a4paper]{amsart}
\usepackage[margin=19mm]{geometry}
\usepackage{amsmath,amssymb,mathtools}
\usepackage[scr=rsfs,cal=euler]{mathalfa}
\usepackage{xfrac}
\usepackage[unicode,hidelinks,bookmarks=false]{hyperref}
\newcommand{\ie}{i.e.\ }
\newcommand{\cf}{cf.\ }
\newcommand{\resp}{respectively\ }
\newcommand{\Subsection}{\subsection}
\providecommand{\nobreakdash}{\nobreak\hbox{-}}
\setlength{\emergencystretch}{2em}
\multlinegap0pt
\usepackage{mathtools}
\usepackage{amssymb}
\usepackage{bm}
\usepackage[shortlabels]{enumitem}
\newtheorem{lemma}{Lemma}
\newcommand{\R}{\mathbb{R}}
\newcommand{\boldlambda}{\boldsymbol{\lambda}}
\newcommand{\calC}{\mathcal C}
\newcommand{\calM}{\mathcal M}
\newcommand{\calX}{\mathcal X}
\newcommand{\cvec}{\mathbf{c}}
\newcommand{\mvec}{\mathbf{m}}
\newcommand{\uvec}{\mathbf{u}}
\newcommand{\vvec}{\mathbf{v}}
\newcommand{\xvec}{\mathbf{x}}
\title{Dual Representation of Minimum Divergence Under Integral Constraints}
\author{Shubhanshu Shekhar, Shubhada Agrawal}
\date{Source: arXiv:2603.21027v1 (CC BY 4.0)}
\begin{document}
\maketitle

\noindent\emph{Source and licence.} Dual Representation of Minimum Divergence Under Integral Constraints. Version arXiv:2603.21027v1 (CC BY 4.0), distributed by arXiv under the Creative Commons Attribution 4.0 International licence, http://creativecommons.org/licenses/by/4.0/, as recorded in the arXiv licence metadata for this exact version and shown on the abstract page https://arxiv.org/abs/2603.21027v1. Full licence notice: \texttt{licenses/arxiv-2603.21027-CC-BY-4.0.txt}.\par\smallskip

\noindent\textit{Editorial scope.} Counted unit: the article's lemma lemma:KL-general-limit-1 (source0.tex lines 1675--1678). The article's complete original proof of it is delivered with it (source0.tex lines 1680--1740, one \texttt{proof} environment of 61 lines). No mathematics is written, repaired or re-derived: the statement and the proof are the article's own lines, in the article's own notation, and no retained line is substituted or re-flowed.  Retained uncounted as the source-local context the proof needs: lines 1667--1669.  Nothing was omitted from any retained span, so the package carries no omission record.  No retained line contains a citation, so the wrapper carries no bibliography and no bibliography entry.  No retained line contains a figure environment, an image-inclusion command or drawing code, so no image is delivered and the package declares no asset dependency.  Editorial formatting only, in addition: an AMS \texttt{amsart} preamble that restates the article's own declarations and macros used by the retained lines, namely the environments \texttt{lemma} on separate counters, together with the standard packages those lines need.  The retained environments print in this document as Lemma 1 in order of appearance, whereas the article prints the same items with the numbering of its own full document.  The complete native source of the article as distributed in the arXiv e-print for this exact version is bundled unchanged as \texttt{sources/196911/source0.tex}.
\begin{align}
\mathbb E_{P_k}\big[\log(\gamma_k^\star-\langle\boldlambda_k^\star,g(X)\rangle)\big] +1-\gamma_k^\star +\inf_{\cvec\in \calC_k}\langle \boldlambda_k^\star,\cvec \rangle \;\ge\;0. \label{eq:opt-nonneg}
\end{align}

\begin{lemma}\label{lemma:KL-general-limit-1}
There exists a constant $R<\infty$ such that for all $k$ large enough, there exists a pair of dual optimizers
$(\boldlambda_k^\star,\gamma_k^\star)$ satisfying $\|(\boldlambda_k^\star,\gamma_k^\star)\|_2\le R$.
\end{lemma}

\begin{proof}
Since $\mvec^{\circ}\in \mathring{\calM}$, where $\calM\coloneqq \mathrm{conv}(g(\calX))$, we have that for every $\uvec \in \R^J$ with $\|\uvec\|_1=1$, $\langle \uvec,\mvec^{\circ}\rangle < \max_{\mvec\in \calM}\langle \uvec,\mvec\rangle$.
The map $\uvec\mapsto \langle \uvec,\mvec^{\circ}\rangle-\max_{\mvec\in \calM}\langle \uvec,\mvec\rangle$ is continuous on the compact
set $\{\uvec:\|\uvec\|_1=1\}$, hence there exists a uniform margin
\begin{align}
\nu \;\coloneqq \;\inf_{\|\uvec\|_1=1}\Big(\max_{\mvec\in \calM}\langle \uvec,\mvec\rangle-\langle \uvec,\mvec^{\circ}\rangle\Big)\;>\;0.
\end{align}

We will now use this ``margin'' property to derive a uniform bound on dual maximizers. Fix a $k\ge 1$, and select any
$(\boldlambda,\gamma)\in \widetilde\Theta_k$. If $r\coloneqq \|\boldlambda\|_1$ is equal to $0$, there is nothing to prove, so assume $r>0$ and set
$\uvec\coloneqq \boldlambda/r$ so that $\|\uvec\|_1=1$. Define the following terms:
\begin{align}
M(\uvec)\coloneqq \max_{\xvec\in\calX}\langle \uvec,g(\xvec)\rangle, \qquad M_k(\uvec)\coloneqq \max_{\vvec\in V_k}\langle \uvec,g(\vvec)\rangle, \qquad \gamma'\coloneqq \gamma-r\,M_k(\uvec).
\end{align}
Feasibility of $(\boldlambda,\gamma)$ implies $\gamma'>0$, and observe that for any $\vvec\in V_k$, we have  $\gamma-\langle \boldlambda,g(\vvec)\rangle
=\gamma' + r\big(M_k(\uvec)-\langle \uvec,g(\vvec)\rangle\big)$.
Since $\|\uvec\|_1=1$ and $\|g(\cdot)\|_\infty\le G_\infty$, this implies $\langle \uvec,g(\vvec)\rangle\in[-G_\infty,G_\infty]$, and hence
\begin{align}
0\le M_k(\uvec)-\langle \uvec,g(\vvec)\rangle\le 2G_\infty \quad\text{and}\quad \gamma-\langle \boldlambda,g(\vvec)\rangle\le \gamma'+2G_\infty r.
\end{align}
Together these bounds imply that $ \mathbb E_{P_k}\!\big[\log(\gamma-\langle \boldlambda,g(X)\rangle)\big]\le \log(\gamma'+2G_\infty r)$. 
Moreover, $\inf_{\cvec\in\calC_k}\langle \boldlambda,\cvec\rangle=r\inf_{\cvec\in\calC_k}\langle \uvec,\cvec\rangle$ and
$\gamma=\gamma'+rM_k(\uvec)$, so the dual objective satisfies
\[
\mathbb E_{P_k}\!\big[\log(\gamma-\langle \boldlambda,g(X)\rangle)\big] + 1-\gamma + \inf_{\cvec\in\calC_k}\langle \boldlambda,\cvec\rangle
\le
\log(\gamma'+2G_\infty r) + 1-\gamma' - r\Big(M_k(\uvec)-\inf_{\cvec\in\calC_k}\langle \uvec,\cvec\rangle\Big).
\label{eq:KL-general-proof-1}
\]

We will now obtain a lower bound on the gap term in \eqref{eq:KL-general-proof-1}. Since $\mvec^{\circ}\in\calC\subseteq \calC_k$, we have $\inf_{\cvec\in\calC_k}\langle \uvec,\cvec\rangle\le \langle \uvec,\mvec^{\circ}\rangle$, hence
$ M_k(\uvec)-\inf_{\cvec\in\calC_k}\langle \uvec,\cvec\rangle \ge M_k(\uvec)-\langle \uvec,\mvec^{\circ}\rangle$. 
Also, because $V_k$ is a $\Delta_k$-cover of $\calX$ and $\eta_k=\omega_g(\Delta_k)$, we have
$ M_k(\uvec)=\max_{\vvec\in V_k}\langle \uvec,g(\vvec)\rangle
\ge \max_{\xvec\in\calX}\langle \uvec,g(\xvec)\rangle-\eta_k
= M(\uvec)-\eta_k$. 
Combining these two inequalities gives us  $M_k(\uvec)-\inf_{\cvec\in\calC_k}\langle \uvec,\cvec\rangle \ge M(\uvec)-\langle \uvec,\mvec^{\circ}\rangle-\eta_k \ge \nu-\eta_k$. Since $\eta_k \downarrow 0$, for all sufficiently large $k$, we have $\eta_k\le \nu/2$, hence $M_k(\uvec)-\inf_{\cvec\in\calC_k}\langle \uvec,\cvec\rangle\ge \nu/2$. Plugging this into \eqref{eq:KL-general-proof-1} yields~(for all sufficiently large $k$),
\begin{align}
\mathbb E_{P_k}\!\big[\log(\gamma-\langle \boldlambda,g(X)\rangle)\big] + 1-\gamma + \inf_{\cvec\in\calC_k}\langle \boldlambda,\cvec\rangle \le \log(\gamma'+2G_\infty r) + 1-\gamma' - (\nu/2)\,r.
\end{align}
Using $\log y\le y-1$ with $y=(\gamma'+2G_\infty r)/(1+2G_\infty r)$ gives
$ \sup_{\gamma'>0}\{\log(\gamma'+2G_\infty r)+1-\gamma'\}\le \log(1+2G_\infty r)+1$, 
and thus for all sufficiently large $k$, we have
\begin{align}
\mathbb E_{P_k}\!\big[\log(\gamma-\langle \boldlambda,g(X)\rangle)\big] + 1-\gamma + \inf_{\cvec\in\calC_k}\langle \boldlambda,\cvec\rangle
\le \log(1+2G_\infty r)+1-(\nu/2)\,r.
\end{align}
The right-hand side tends to $-\infty$ as $r\to\infty$. Since any maximizer $(\boldlambda_k^\star,\gamma_k^\star)$ satisfies \eqref{eq:opt-nonneg}, it follows that $\|\boldlambda_k^\star\|_1$ is uniformly bounded for all sufficiently large $k$. Let us call the bound $R_{\boldlambda}$.

Finally, feasibility gives $\gamma_k^\star>\max_{\vvec\in V_k}\langle \boldlambda_k^\star,g(\vvec)\rangle\ge
-\|\boldlambda_k^\star\|_1 G_\infty\ge -R_{\boldlambda}G_\infty$, giving a uniform lower bound on $\gamma_k^\star$.
Moreover, since $\calC$ is compact and $\eta_k\to 0$, the sets $\calC_k$ are uniformly bounded; thus
$\inf_{\cvec\in\calC_k}\langle \boldlambda_k^\star,\cvec\rangle\le \|\boldlambda_k^\star\|_1\sup_{\cvec\in\calC_k}\|\cvec\|_\infty$
is uniformly bounded above. Using again that the term $1-\gamma$ drives the objective to $-\infty$ as $\gamma\to\infty$,
while \eqref{eq:opt-nonneg} must hold at a maximizer, we obtain a uniform upper bound $\gamma_k^\star\le R_\gamma$ for all sufficiently
large $k$. Enlarging constants if needed, there exists $R<\infty$ such that for all sufficiently large $k$,
\[
\|(\boldlambda_k^\star,\gamma_k^\star)\|_2\le R. \label{eq:theta-star-bound}
\]
This completes the proof.
\end{proof}

\end{document}
